% ---
% school: "college"
% level: "8eme"
% section: null
% chapter: "العمليات والحساب على الأعداد الكسرية"
% source: "Axiomaths"
% title: "الجمع والطرح"
% country: "Tunisie"
% year: 2026
% professor: "Axiomaths"
% difficulty: "Facile"
% points: null
% tags: ["fractions", "addition", "soustraction"]
% ---
% Exercice 3/15 de la série 8ème « العمليات والحساب على الأعداد الكسرية » (statement + \axsol + solution). Compiler avec XeLaTeX + axiomath.sty (+ axiomaths-logo.png).
\documentclass[12pt,a4paper]{article}
\usepackage[level=8b,color=15803D]{axiomath}
\axsetup{title={العمليات والحساب على الأعداد الكسرية}, corrige}

\begin{document}

\begin{corr}[title={الجمع والطرح}, stars=1, source={Axiomaths}]
احسب ثم اختزل النتيجة إن أمكن :
\begin{enumerate}
  \item $\dfrac{2}{5}+\dfrac{1}{3}$
  \item $\dfrac{5}{6}-\dfrac{3}{4}$
  \item $-\dfrac{3}{10}-\dfrac{4}{15}$
  \item $\dfrac{7}{12}-\left(-\dfrac{5}{18}\right)$
  \item $3-\dfrac{5}{7}$
  \item $-\dfrac{1}{2}-\dfrac{1}{3}-\dfrac{1}{6}$
\end{enumerate}
\axsol
\begin{prop}[قاعدة]
لجمع أو طرح عددين كسريين نوحّد المقامين ثم نجمع أو نطرح البسطين ونحافظ على المقام :
\ib{$\dfrac{a}{c}\pm\dfrac{b}{c}=\dfrac{a\pm b}{c}$}
\end{prop}
\begin{prop}[طرح عدد]
طرح عدد يعني جمع مضاده : \ib{$x-y=x+(-y)$}
\end{prop}
\begin{enumerate}[itemsep=14pt]
  \item المقام المشترك هو $15$ :
  \[\dfrac{2}{5}+\dfrac{1}{3}=\dfrac{6}{15}+\dfrac{5}{15}=\dfrac{11}{15}\]
  \item المقام المشترك هو $12$ :
  \[\dfrac{5}{6}-\dfrac{3}{4}=\dfrac{10}{12}-\dfrac{9}{12}=\dfrac{1}{12}\]
  \item المقام المشترك هو $30$ :
  \[-\dfrac{3}{10}-\dfrac{4}{15}=-\dfrac{9}{30}-\dfrac{8}{30}=-\dfrac{17}{30}\]
  \item طرح $-\dfrac{5}{18}$ يعني جمع $\dfrac{5}{18}$ ، والمقام المشترك هو $36$ :
  \[\dfrac{7}{12}-\left(-\dfrac{5}{18}\right)=\dfrac{7}{12}+\dfrac{5}{18}=\dfrac{21}{36}+\dfrac{10}{36}=\dfrac{31}{36}\]
  \item نكتب $3=\dfrac{21}{7}$ :
  \[3-\dfrac{5}{7}=\dfrac{21}{7}-\dfrac{5}{7}=\dfrac{16}{7}\]
  \item المقام المشترك هو $6$ :
  \[-\dfrac{1}{2}-\dfrac{1}{3}-\dfrac{1}{6}=-\dfrac{3}{6}-\dfrac{2}{6}-\dfrac{1}{6}=-\dfrac{6}{6}=-1\]
\end{enumerate}
\end{corr}

\end{document}
